\documentclass[11pt,a4paper]{article}
\usepackage[margin=0.93in]{geometry}
\usepackage{graphicx,booktabs,multirow}
\usepackage[authoryear]{natbib}
\usepackage{enumitem}
\usepackage[hidelinks]{hyperref}

\setlength{\parindent}{0pt}
\setlength{\parskip}{4pt}
\setlength{\intextsep}{7pt}
\setlist[itemize]{leftmargin=1.6em,topsep=2pt,itemsep=2pt,parsep=0pt}
\renewcommand{\thefigure}{A\arabic{figure}}
\renewcommand{\thetable}{A\arabic{table}}
\newcommand{\arq}[1]{\par\bigskip{\large\bfseries\raggedright #1\par}\smallskip}

\begin{document}
\section*{Appendix Research Questions}

\arq{ARQ1: How effective is LAMIC with different LLMs?}

\noindent\textbf{Motivation:}
In the main experiments, LAMIC uses DeepSeek-V4 to locate relevant
\(\langle\mathrm{API},\mathrm{KU}\rangle\) pairs. This raises the
question of whether LAMIC remains effective when it uses a different
LLM. In particular, smaller open-source LLMs can be deployed locally,
making them a practical choice for this task. This ARQ aims to evaluate
how the choice of LLM affects the effectiveness of LAMIC on the Java
and Android datasets.

\noindent\textbf{Baselines:}
We compare LAMIC with three variants to locate relevant
\(\langle\mathrm{API},\mathrm{KU}\rangle\) pairs. All four settings use the same demonstration clues generated by DeepSeek-V4.

\begin{itemize}
    \item \textbf{\(\mathrm{Model}_{\mathrm{Llama}}\)}~\citep{MetaLlama31}
    uses Llama-3.1-8B, deployed locally via
    Ollama\footnote{\url{https://ollama.com/library/llama3.1:8b}}
    on a single NVIDIA A800 GPU.

    \item \textbf{\(\mathrm{Model}_{\mathrm{Qwen}}\)}~\citep{qwen38}
    uses Qwen-3.8-27B from the Ollama model
    library\footnote{\url{https://ollama.com/library/qwen3.8:27b}},
    deployed on the same A800 GPU.

    \item \textbf{\(\mathrm{Model}_{\mathrm{GPT}}\)}~\citep{openai2026terra}
    uses GPT-6.1-Sol, accessed through the OpenAI
    API\footnote{\url{https://developers.openai.com/api/docs/models/gpt-6.1-sol}}.
\end{itemize}

\begin{figure}[!htbp]
    \centering
    \includegraphics[width=0.78\linewidth,height=0.24\textheight,keepaspectratio]{ARQ1_cropped (2).pdf}
    \caption{Average Results of LAMIC and the Baselines Using
    Different LLMs (ARQ1)}
    \label{fig:appendix-arq1}
\end{figure}

\noindent\textbf{Results:}
As shown in Figure~\ref{fig:appendix-arq1}, averaged across the Java and
Android datasets, \(\mathrm{Model}_{\mathrm{Llama}}\) and
\(\mathrm{Model}_{\mathrm{Qwen}}\) achieve F1 scores of 82.61\% and
83.54\%, respectively, compared with 86.12\% for LAMIC.
\(\mathrm{Model}_{\mathrm{GPT}}\) achieves the highest average F1 of
87.83\%, outperforming LAMIC by 1.71 percentage points. This may be
because GPT-6.1-Sol is a recent LLM with stronger capabilities.
The improvement is statistically significant with \(p=0.0187\) and
has a small effect size with Cliff's \(\delta=0.2565\).

Although \(\mathrm{Model}_{\mathrm{GPT}}\) outperforms LAMIC in F1,
LAMIC has a lower cost per output token. Regarding dollar cost,
GPT-6.1-Sol costs \$10.00 per million output tokens, and
DeepSeek-V4 costs \$3.96 per million output tokens. Therefore,
we decided to use DeepSeek-V4 as the backbone for the main experiments.

\noindent\textbf{In summary,} LAMIC can generalize to different LLMs
when locating relevant \(\langle\mathrm{API},\mathrm{KU}\rangle\) pairs.

\clearpage
\arq{ARQ2: How effective is Standard ICL in locating relevant API knowledge?}

\noindent\textbf{Motivation:}
To the best of our knowledge, LAMIC is the first approach to use ICL
to locate relevant \(\langle\mathrm{API},\mathrm{KU}\rangle\) pairs.
LAMIC extends Standard ICL with two strategies, namely the
multi-perspective demonstration selection strategy and the clue-based
demonstration enhancement strategy. Standard ICL can also be applied
directly to this task. This ARQ aims to evaluate how effectively
Standard ICL locates relevant \(\langle\mathrm{API},\mathrm{KU}\rangle\)
pairs on the Java and Android datasets.

\noindent\textbf{Baseline:}
We compare LAMIC with the following baseline. Both approaches use DeepSeek-V4 as the underlying LLM.

\begin{itemize}
    \item \textbf{Standard ICL}~\citep{brown2020language} is a few-shot learning approach that leverages the ability of LLMs to learn from demonstrations provided in the input prompt. Specifically, we use SentenceBERT to retrieve five semantically similar $\langle$API, KU$\rangle$ pairs as demonstrations, and employ DeepSeek-V4 for locating relevant $\langle$API, KU$\rangle$ pairs.
\end{itemize}

\begin{table}[!htbp]
\centering
\caption{Average Results of LAMIC and Baseline Approaches in ARQ2 (\%)}
\label{tab:appendix-arq2}
\small
\setlength{\tabcolsep}{0pt}
\begin{tabular*}{\columnwidth}{@{\extracolsep{\fill}}lcccccc@{}}
\toprule
\multirow{2}{*}{Approach}
  & \multicolumn{3}{c}{Java Dataset}
  & \multicolumn{3}{c}{Android Dataset} \\
\cmidrule(lr){2-4}\cmidrule(l){5-7}
  & P & R & F1 & P & R & F1 \\
\midrule
GMR          & 49.10 & 46.78 & 46.15 & 49.54 & 48.55 & 45.93 \\
FITSEA       & 56.60 & 55.55 & 51.89 & 53.30 & 55.54 & 51.33 \\
FRAPT        & 51.70 & 69.00 & 58.92 & 44.44 & 65.71 & 52.98 \\
Standard ICL & 67.65 & 77.08 & 70.26 & 53.67 & 69.97 & 60.43 \\
UniXcoder    & 74.89 & 76.09 & 73.55 & 73.16 & 66.72 & 68.35 \\
DTML         & 80.73 & 76.26 & 76.13 & 86.71 & 81.47 & 80.26 \\
\midrule
LAMIC        & \textbf{81.81} & \textbf{94.29} & \textbf{87.41}
             & \textbf{88.12} & \textbf{82.85} & \textbf{84.82} \\
\bottomrule
\end{tabular*}
\end{table}

\noindent\textbf{Results:}
As shown in Table~\ref{tab:appendix-arq2}, LAMIC outperforms Standard ICL
by an average of 20.77 percentage points in F1 across the Java and
Android datasets. This improvement may be attributed to two strategies
in LAMIC. The multi-perspective demonstration selection strategy fuses
the lexical, semantic, and structural rankings to retrieve candidate
demonstrations for ICL. The clue-based demonstration enhancement
strategy augments each retrieved candidate demonstration with
explicit clues and a reasoning statement.

\noindent\textbf{In summary,} our proposed framework LAMIC improves
Standard ICL with a multi-perspective demonstration selection strategy
and a clue-based demonstration enhancement strategy, enabling it to
locate relevant \(\langle\mathrm{API},\mathrm{KU}\rangle\) pairs more
effectively.

\clearpage
\arq{ARQ3: How effective is LAMIC with different semantic retrievers?}

\noindent\textbf{Motivation:}
LAMIC uses SentenceBERT~\citep{reimers2019sentencebert} to retrieve
semantically similar demonstrations. Since KUs can contain both natural
language and code, a code-aware model could also serve as the semantic
retriever. This ARQ aims to evaluate how the choice of semantic retriever
affects the effectiveness of LAMIC in locating relevant
\(\langle\mathrm{API},\mathrm{KU}\rangle\) pairs.

\noindent\textbf{Baselines:}
We compare two variants of LAMIC that use different semantic retrievers
for demonstration selection.

\begin{itemize}
    \item \textbf{\(\mathrm{Model}_{\mathrm{SB}}\)}~\citep{reimers2019sentencebert}
    uses SentenceBERT to encode each test pair and the candidate
    demonstrations and ranks the candidates by cosine similarity.
    This is the semantic retriever used in LAMIC.

    \item \textbf{\(\mathrm{Model}_{\mathrm{CB}}\)}~\citep{CodeBERT}
    uses CodeBERT to encode each test pair and the candidate
    demonstrations. It computes the cosine similarity between their
    embeddings to obtain the semantic ranking.
\end{itemize}

\begin{figure}[!htbp]
    \centering
    \includegraphics[width=0.78\linewidth,height=0.24\textheight,keepaspectratio]{ARQ3.pdf}
    \caption{Average Results of Different Semantic Retrievers (ARQ3)}
    \label{fig:appendix-arq3}
\end{figure}

\noindent\textbf{Results:}
As shown in Figure~\ref{fig:appendix-arq3}, LAMIC outperforms
\(\mathrm{Model}_{\mathrm{CB}}\) in F1 by 3.35 percentage points on
Java and 1.66 percentage points on Android. One possible reason is
that tutorial segments and SO answers often explain in natural
language how an API is used. SentenceBERT encodes each passage as a
whole, which may help LAMIC retrieve demonstrations that describe
similar API uses.

\noindent\textbf{In summary,}
LAMIC achieves higher F1 with SentenceBERT than with CodeBERT on both
datasets. We therefore retain SentenceBERT for semantic demonstration
retrieval.

\clearpage
\arq{ARQ4: How effective is LAMIC on imbalanced SO datasets?}

\noindent\textbf{Motivation:}
We evaluated LAMIC on balanced SO datasets with equal numbers of
relevant and irrelevant \(\langle\mathrm{API},\mathrm{QA}\rangle\)
pairs. This ARQ aims to evaluate how precisely LAMIC locates relevant
pairs on imbalanced SO datasets, where irrelevant pairs outnumber
relevant ones.

\noindent\textbf{Baselines:}
We compare the balanced and imbalanced settings for the same five
Java and four Android libraries. The balanced setting combines API
tutorials with balanced SO subsets. The imbalanced setting retains
the class distribution of each SO dataset. Across the five Java
libraries, the imbalanced SO datasets contain 13,395
\(\langle\mathrm{API},\mathrm{QA}\rangle\) pairs, including
3,608 relevant and 9,787 irrelevant pairs. Across the four Android
libraries, the corresponding totals are 7,613 pairs, including
1,309 relevant and 6,304 irrelevant pairs.

\begin{figure}[!htbp]
    \centering
    \includegraphics[width=0.78\linewidth,height=0.24\textheight,keepaspectratio]{R3Q2_cropped.pdf}
    \caption{Average Results of LAMIC under the Balanced and
    Imbalanced Settings (ARQ4)}
    \label{fig:appendix-arq4}
\end{figure}

\noindent\textbf{Results:}
As shown in Figure~\ref{fig:appendix-arq4}, LAMIC achieves an average F1
of 81.75\% on Java and 84.27\% on Android under the imbalanced
setting. Under the balanced setting, average F1 increases to
87.41\% on Java and 84.82\% on Android. A possible reason is that
ICL is sensitive to the label distribution of its
demonstrations~\citep{Yang2024Harnessing,Song2024Shortcut}.

\noindent\textbf{In summary,} when evaluated on imbalanced SO datasets,
LAMIC achieves an average F1 of 81.75\% on Java and 84.27\% on Android.

\clearpage
\bibliographystyle{elsarticle-harv}
\bibliography{references}
\end{document}
